\hspace{3mm}

\centerline{\bf \Huge Домашнее задание.}
\centerline{\bf \Huge Турниры}

\begin{flushright}
    \scriptsize{Насладись небольшим отклонением от пути.\\ По полной.\\ Там ты сможешь найти то, что важнее желаемого тобой. \\Джин Фрикс. Хантер}
\end{flushright}

\problem{1!} Весной в ЭлМШ будет проведен турнир X по шахматам среди учеников, в котором каждый сыграет с каждым из игроков один раз. Определите общее количество партий, которое будет сыграно, если известно, что в ЭлМШ обучаются 50 учеников? А сколько партий сыграет один ученик?

\problem{2!} Какое суммарное количество очков наберут все участника турнира L из задачи 1!. , если за победу в партии присуждается 2 очка, за ничью - 1, а за поражение - 0?

\problem{3} В пятикруговом турнире по футболу принимают участие 100 команд. Определите количество матчей, которые были проведены за весь турнир? А сколько матчей проведет одна команда?

\problem{4} В однокруговом турнире по футболу соревнуются 20 команд. За победу присуждается 3 очка, за ничью - 1, за поражение - 0. Известно, что по окончанию турнира Реал Мадрид набрал 57 очков. Могла ли Барселона набрать 55 очков?